\section{Related Work}
\label{related_work}

Previous work in model-based reinforcement learning has focused on planning in
low-dimensional state spaces \linebreak \citep{gal2016deeppilco,higuera2018synthesizing,henaff2018planbybackprop,chua2018pets}, combining the benefits of model-based and model-free approaches \citep{kalweit2017blending,nagabandi2017mbmf,weber2017i2a,kurutach2018modeltrpo,buckman2018steve,ha2018worldmodels,wayne2018merlin,igl2018dvrl,srinivas2018upn}, and pure video prediction without planning \citep{oh2015atari,krishnan2015deepkalman,karl2016dvbf,chiappa2017recurrent,babaeizadeh2017sv2p,gemici2017temporalmemory,denton2018stochastic,buesing2018dssm,doerr2018prssm,gregor2018tdvae}. \Cref{sec:more_related_work} reviews these orthogonal research directions in more detail.

Relatively few works have demonstrated successful planning from pixels using learned dynamics models. The robotics community focuses on video prediction models for planning \citep{agrawal2016poking,finn2017foresight,ebert2018foresight,zhang2018solar} that deal with the visual complexity of the real world and solve tasks with a simple gripper, such as grasping or pushing objects. In comparison, we focus on simulated environments, where we leverage latent planning to scale to larger state and action spaces, longer planning horizons, as well as sparse reward tasks. E2C \citep{watter2015e2c} and RCE \citep{banijamali2017rce} embed images into a latent space, where they learn local-linear latent transitions and plan for actions using LQR. These methods balance simulated cartpoles and control 2-link arms from images, but have been difficult to scale up. We lift the Markov assumption of these models, making our method applicable under partial observability, and present results on more challenging environments that include longer planning horizons, contact dynamics, and sparse rewards.
